% Generated draft: requires mathematical and source review.
\documentclass[11pt]{article}
\usepackage[margin=25mm]{geometry}
\usepackage{fontspec}
\setmainfont{Latin Modern Roman}
\usepackage{amsmath,amssymb,mathtools,unicode-math}
\setmathfont{Latin Modern Math}
\usepackage{enumitem,needspace,xurl,hyperref}
\setlength{\parindent}{0pt}
\setlength{\parskip}{5pt}
\setlength{\emergencystretch}{3em}
\newcommand{\N}{\mathbb{N}}
\newcommand{\Z}{\mathbb{Z}}
\newcommand{\Q}{\mathbb{Q}}
\newcommand{\R}{\mathbb{R}}
\newcommand{\C}{\mathbb{C}}
\newcommand{\F}{\mathbb{F}}
\newcommand{\A}{\mathbb{A}}
\newcommand{\PP}{\mathbb{P}}
\newcommand{\E}{\mathbb{E}}
\newcommand{\eps}{\varepsilon}
\newcommand{\dd}{\,\mathrm{d}}
\begin{document}
\section*{UnsolvedMath ID: 30000642}
\textbf{Source problem code:} catalogue entry
\subsection{Definitions and mathematical statement}
Let \(d\geq3\), let \(R\) be the characteristic-zero ring intended in the source, and put \(A=R[x_1,\ldots,x_d]\). In the usual commutative unital convention, characteristic zero does not imply that \(R\) is a field, a domain, or a \(\Q\)-algebra, nor that its residue fields have characteristic zero. No such stronger assumption is supplied.

A coordinate of \(A\) over \(R\) is an element \(f\) extendable to polynomial coordinates, so that \(A=R[f,g_2,\ldots,g_d]\) as a polynomial algebra. A usual residual-variable convention requires that for every prime \(\mathfrak p\subset R\), the image \(\overline f\) be a coordinate of \(\kappa(\mathfrak p)[x_1,\ldots,x_d]\). A ``bad residual variable'' might mean one satisfying this fiberwise property but not the global coordinate property; the source does not reproduce its definition and that interpretation must be checked. Merely requiring each fiber quotient by \(\overline f\) to be affine space is a different condition and is not silently identified with fiberwise coordinateness.

Similarly, an algebraic hyperplane may mean a hypersurface \(V(f)\subset\operatorname{Spec}A\) isomorphic over \(R\) to \(\A_R^{d-1}\), equivalently \(A/(f)\cong R^{[d-1]}\). Calling it bad could refer to failure of rectification to a coordinate hyperplane, noncoordinateness of \(f\), or a stronger source-specific property. It need not mean a linear degree-one hyperplane. The corresponding bad-line constructions and their precise defect are not supplied.

The source task is to generalize those constructions to obtain bad hyperplanes and bad residual variables in \(d\)-variable polynomial rings over the intended characteristic-zero bases. It does not give a construction to extend, define bad, state whether every such \(R\) is intended or only suitable examples, or prescribe the additional properties to be preserved. Therefore neither a universal existence theorem for all characteristic-zero rings nor an invented higher-dimensional example is asserted.

The supplied literature triage checked 2026-08-21 reports no authoritative complete resolution of the extracted generalization request, without separate definitions or theorem citations.
\subsection{Short English statement}
Extend the source's bad-line and residual-variable constructions to \(d\geq3\). The meaning of bad, the original construction, and the allowed base rings are missing; fiberwise coordinates, affine-space quotients, and global rectifiability must remain distinct.
\subsection{Sources}
\begin{itemize}
\item Local UnsolvedMath registry, permanent ID 30000642, catalogue entry, \emph{Bad Hyperplanes and Residual Variables in Polynomial Rings}.
\item \emph{Affine Algebraic Geometry} (2007), Oberwolfach Reports, pp. 5--82, DOI \url{https://doi.org/10.4171/owr/2007/01}.
\item Supplied literature triage checked 2026-08-21 cites only the original report; no bad-line construction, formal residual-variable definition, or separate theorem reference is supplied.
\end{itemize}
\end{document}
